\chapter{3D 并行策略}
\label{chap:3d-parallelism}

\section{DP、TP、PP 的组合}
\label{sec:dp-tp-pp-composition}

前面三章分别讨论了数据并行、张量并行和流水线并行。到这一章，我们把三者组合起来，进入大模型训练中最常见的系统形态：\textbf{3D 并行}。

所谓 3D 并行，并不是指某个单独的新算法，而是把三类并行维度组织成一个三维并行网格：

[
\text{Data Parallelism}
\times
\text{Tensor Parallelism}
\times
\text{Pipeline Parallelism}.
]

如果用符号表示：

[
W
=

P_{\mathrm{DP}}
\cdot
P_{\mathrm{TP}}
\cdot
P_{\mathrm{PP}},
]

其中：

\begin{itemize}
\item $W$ 是总 GPU 数，也就是 world size；
\item $P_{\mathrm{DP}}$ 是数据并行 size；
\item $P_{\mathrm{TP}}$ 是张量并行 size；
\item $P_{\mathrm{PP}}$ 是流水线并行 size。
\end{itemize}

3D 并行的目标是把一个超大 Transformer 训练任务拆解到许多 GPU 上，同时控制三个核心成本：

\begin{itemize}
\item \textbf{显存}：模型参数、梯度、优化器状态、激活和通信 buffer 能否放下；
\item \textbf{计算}：每张 GPU 是否有足够大的 GEMM 和稳定工作量；
\item \textbf{通信}：DP、TP、PP 三类通信是否能被网络拓扑和调度承受。
\end{itemize}

如果只使用一种并行方式，往往会遇到边界。只用 DDP，单卡必须放下完整模型；只用 TP，层内通信过于频繁，TP size 不能无限增大；只用 PP，pipeline bubble 和 stage 负载均衡会限制效率。因此，大模型训练通常需要三者组合。

\subsection{三类并行的正交性}
\label{subsec:orthogonality-of-parallelism}

DP、TP、PP 的“正交性”是理解 3D 并行的关键。

\textbf{数据并行}切分 batch。每个数据并行副本看到不同数据分片，但副本之间逻辑上训练同一个模型。每个 DP replica 内部可以包含多个 TP ranks 和多个 PP stages。DP 维度的通信主要是梯度同步，常见操作包括 all-reduce、reduce-scatter 或 ZeRO/FSDP 中的参数/梯度分片通信。

[
\text{DP}: \quad
\mathcal{B}
===========

\mathcal{B}*0
\cup
\mathcal{B}*1
\cup
\cdots
\cup
\mathcal{B}*{P*{\mathrm{DP}}-1}.
]

\textbf{张量并行}切分每一层内部的大矩阵。一个 Transformer block 内的 QKV projection、output projection、FFN up/gate/down projection 会被切到 TP group 内的多张 GPU 上。TP 维度通信发生在每层内部，频率很高，因此通常要求 TP group 位于同一节点内或高速互联域内。

[
\mat{W}
=======

[\mat{W}*0,\mat{W}*1,\ldots,\mat{W}*{P*{\mathrm{TP}}-1}]
\quad
\text{or}
\quad
\mat{W}
=======

\begin{bmatrix}
\mat{W}*0\
\mat{W}*1\
\vdots\
\mat{W}*{P*{\mathrm{TP}}-1}
\end{bmatrix}.
]

\textbf{流水线并行}切分模型层。不同 PP stages 负责不同层段，micro-batches 像流水线一样依次穿过这些 stages。PP 维度通信主要是相邻 stage 之间发送 forward activation 和 backward activation gradient。

[
{0,1,\ldots,L-1}
================

\mathcal{L}*0
\cup
\mathcal{L}*1
\cup
\cdots
\cup
\mathcal{L}*{P*{\mathrm{PP}}-1}.
]

三者正交的意思是：同一个全局 rank 可以同时属于一个 DP group、一个 TP group 和一个 PP group。它在 DP 维度上代表某个数据副本；在 TP 维度上代表某个 tensor shard；在 PP 维度上代表某个 layer stage。

可以把一个 rank 的身份写成三元组：

[
r
\leftrightarrow
(d,t,p),
]

其中：

\begin{itemize}
\item $d\in[0,P_{\mathrm{DP}})$ 表示 data parallel index；
\item $t\in[0,P_{\mathrm{TP}})$ 表示 tensor parallel index；
\item $p\in[0,P_{\mathrm{PP}})$ 表示 pipeline parallel index。
\end{itemize}

一个完整的全局 rank 映射函数可以写为：

[
r
=

((d\cdot P_{\mathrm{PP}})+p)\cdot P_{\mathrm{TP}}+t.
]

也可以使用其他排列方式，例如让 TP 维最快变化，PP 维次之，DP 维最慢变化。不同排列方式会影响 rank mapping 和物理拓扑匹配，这一点稍后会详细讨论。

\begin{figure}[H]
\centering
\begin{tikzpicture}[
font=\small,
cube/.style={draw, rounded corners=3pt, minimum width=1.1cm, minimum height=0.65cm, align=center, fill=blue!6},
arrow/.style={-{Latex[length=2.2mm]}, thick},
group/.style={draw, rounded corners=4pt, dashed, thick}
]

```
\node[font=\bfseries] at (0,4.0) {3D 并行：DP、TP、PP 三个正交维度};

% Draw a 2 x 2 x 3-ish grid projected
\foreach \d/\dx in {0/0,1/0.55} {
  \foreach \p/\py in {0/0,1/0.9,2/1.8} {
    \foreach \t/\tx in {0/0,1/1.25} {
      \pgfmathsetmacro{\x}{-3.6+\tx+\dx}
      \pgfmathsetmacro{\y}{-0.5+\py+\dx*0.35}
      \node[cube] (r\d\t\p) at (\x,\y) {$d\d$\\$t\t,p\p$};
    }
  }
}

\draw[arrow] (-4.4,-1.2) -- (-2.3,-1.2) node[midway,below] {TP};
\draw[arrow] (-4.55,-0.85) -- (-4.55,2.1) node[midway,left] {PP};
\draw[arrow] (-4.25,-0.95) -- (-3.55,-0.55) node[midway,below right] {DP};

\node[group, fit=(r000)(r010)] {};
\node[align=center, font=\scriptsize] at (-2.95,0.1) {同一层段内\\TP group};

\draw[green!60!black, thick, dotted] (r000) -- (r100);
\draw[green!60!black, thick, dotted] (r010) -- (r110);
\node[align=center, font=\scriptsize, green!50!black] at (-1.2,-0.25) {DP replicas};

\draw[red!70, thick, dashed] (r000) -- (r001) -- (r002);
\node[align=center, font=\scriptsize, red!70] at (-4.85,1.4) {PP stages};

\node[align=left, text width=10.8cm] at (1.8,-2.0) {
  每个 rank 同时有 DP、TP、PP 三个坐标。TP group 负责层内张量切分；
  PP group 负责层间流水线；DP group 负责不同数据副本之间的梯度同步。
};
```

\end{tikzpicture}
\caption{3D 并行 rank 拓扑示意图}
\label{fig:3d-parallel-rank-topology}
\end{figure}

\subsection{group 划分}
\label{subsec:3d-parallel-group-partition}

在实际分布式训练框架中，3D 并行不是抽象概念，而是具体的 process groups。每个 rank 需要加入多个通信组。最少包括：

\begin{itemize}
\item \textbf{Tensor Parallel Group}：组内做层内 all-reduce、all-gather、reduce-scatter；
\item \textbf{Pipeline Parallel Group}：组内相邻 rank 做 send/recv；
\item \textbf{Data Parallel Group}：组内同步相同模型 shard 的梯度或状态；
\item \textbf{Embedding Group}：首尾 stage 之间共享或同步 embedding/lm head 时可能需要；
\item \textbf{Sequence Parallel Group}：通常与 TP group 相同或相关；
\item \textbf{Expert Parallel Group}：MoE 训练中用于 expert dispatch 和 all-to-all。
\end{itemize}

假设：

[
P_{\mathrm{TP}}=2,\quad
P_{\mathrm{PP}}=4,\quad
P_{\mathrm{DP}}=2.
]

则总 GPU 数为：

[
W=2\times 4\times 2=16.
]

如果使用 rank 映射：

[
r=((dP_{\mathrm{PP}})+p)P_{\mathrm{TP}}+t,
]

则全局 rank 表如下：

[
\begin{array}{c|c|c|c}
d & p & t & r\
\hline
0 & 0 & 0 & 0\
0 & 0 & 1 & 1\
0 & 1 & 0 & 2\
0 & 1 & 1 & 3\
0 & 2 & 0 & 4\
0 & 2 & 1 & 5\
0 & 3 & 0 & 6\
0 & 3 & 1 & 7\
1 & 0 & 0 & 8\
1 & 0 & 1 & 9\
1 & 1 & 0 & 10\
1 & 1 & 1 & 11\
1 & 2 & 0 & 12\
1 & 2 & 1 & 13\
1 & 3 & 0 & 14\
1 & 3 & 1 & 15
\end{array}
]

此时 TP groups 为：

[
(0,1),(2,3),(4,5),(6,7),
]
[
(8,9),(10,11),(12,13),(14,15).
]

PP groups 为：

[
(0,2,4,6),\quad (1,3,5,7),
]
[
(8,10,12,14),\quad (9,11,13,15).
]

DP groups 为：

[
(0,8),(1,9),(2,10),(3,11),
]
[
(4,12),(5,13),(6,14),(7,15).
]

注意 DP group 中的 rank 必须代表\textbf{相同的 TP shard 和相同的 PP stage}，只是数据副本不同。因此 rank $0$ 和 rank $8$ 都是 $t=0,p=0$，只是 $d$ 不同。它们之间同步的是同一 shard 的梯度。

下面给出一个简单的 group 生成代码。它用于帮助理解，真实框架会考虑更多并行维度和拓扑策略。

\begin{lstlisting}[language=Python,caption={3D 并行中 TP/PP/DP groups 的生成示意},label={lst:3d-parallel-group-generation}]
def rank_from_coords(d, p, t, dp, pp, tp):
# Layout: DP outer, PP middle, TP inner.
return ((d * pp) + p) * tp + t

def build_3d_parallel_groups(dp, pp, tp):
world_size = dp * pp * tp

```
tp_groups = []
for d in range(dp):
    for p in range(pp):
        group = [
            rank_from_coords(d, p, t, dp, pp, tp)
            for t in range(tp)
        ]
        tp_groups.append(group)

pp_groups = []
for d in range(dp):
    for t in range(tp):
        group = [
            rank_from_coords(d, p, t, dp, pp, tp)
            for p in range(pp)
        ]
        pp_groups.append(group)

dp_groups = []
for p in range(pp):
    for t in range(tp):
        group = [
            rank_from_coords(d, p, t, dp, pp, tp)
            for d in range(dp)
        ]
        dp_groups.append(group)

assert sorted(sum(tp_groups, [])) == list(range(world_size))
return tp_groups, pp_groups, dp_groups
```

tp_groups, pp_groups, dp_groups = build_3d_parallel_groups(
dp=2,
pp=4,
tp=2,
)

print("TP groups:", tp_groups)
print("PP groups:", pp_groups)
print("DP groups:", dp_groups)
\end{lstlisting}

在 PyTorch 中，这些 group 通常会通过 \texttt{torch.distributed.new_group()} 创建。训练框架会保存每个 rank 所属 group，并在不同模块中调用对应 group 的 collective 或 point-to-point 通信。

\subsection{rank mapping}
\label{subsec:rank-mapping}

Rank mapping 是把逻辑三维坐标 $(d,p,t)$ 映射到物理 GPU 的过程。它极其重要，因为不同并行维度的通信频率和通信模式不同。

一般经验是：

\begin{itemize}
\item \textbf{TP group 优先放在节点内高速互联}，因为 TP 通信发生在每层内部，频率最高；
\item \textbf{PP 相邻 stage 尽量放在高速链路上}，因为每个 micro-batch 都要传递激活和梯度；
\item \textbf{DP 可以跨节点}，因为 DP 通信虽然数据量大，但通常可与 backward 重叠，且模式规则；
\item \textbf{同一节点内的 GPU 拓扑要考虑 NVLink/NVSwitch/PCIe 层级}；
\item \textbf{跨节点要考虑 NIC 数量、rail、RDMA 路径和机架拓扑}。
\end{itemize}

设一个节点有 $G_{\mathrm{node}}=8$ 张 GPU。如果选择：

[
P_{\mathrm{TP}}=8,
]

则一个 TP group 可以完整放在单节点内。这样每层 TP all-reduce 都走节点内 NVLink/NVSwitch，通常性能较好。如果选择：

[
P_{\mathrm{TP}}=16,
]

则一个 TP group 跨两个节点，每层 TP 通信都要走跨节点网络，性能可能明显下降。除非模型极大必须这么做，否则通常避免跨节点 TP。

PP group 是否跨节点要看 stage 数和节点布局。如果每个 stage 是一个 TP group，而 TP group 占满一个节点，那么 PP 相邻 stage 自然跨节点。此时 PP 通信是激活传递，其大小与：

[
B_{\mathrm{micro}}S H
]

成正比。对于长序列和大 hidden，PP 通信也不可忽略。因此有些训练会把相邻 PP stages 安排在网络拓扑相近的节点上。

\begin{figure}[H]
\centering
\begin{tikzpicture}[
font=\small,
nodebox/.style={draw, rounded corners=5pt, minimum width=5.0cm, minimum height=2.3cm, fill=blue!3},
gpu/.style={draw, rounded corners=2pt, minimum width=0.7cm, minimum height=0.46cm, align=center, fill=orange!12},
arrow/.style={-{Latex[length=2.0mm]}, thick},
fast/.style={thick, green!60!black},
slow/.style={thick, red!70, dashed}
]

```
\node[font=\bfseries] at (0,3.5) {Rank Mapping：让高频通信走高速链路};

\node[nodebox, label={[font=\bfseries]above:Node 0}] (n0) at (-3.1,1.1) {};
\node[nodebox, label={[font=\bfseries]above:Node 1}] (n1) at (3.1,1.1) {};

\foreach \i/\x/\y in {0/-4.6/1.55,1/-3.8/1.55,2/-3.0/1.55,3/-2.2/1.55,4/-4.6/0.75,5/-3.8/0.75,6/-3.0/0.75,7/-2.2/0.75} {
  \node[gpu] (g\i) at (\x,\y) {\i};
}

\foreach \i/\x/\y/\r in {0/1.6/1.55/8,1/2.4/1.55/9,2/3.2/1.55/10,3/4.0/1.55/11,4/1.6/0.75/12,5/2.4/0.75/13,6/3.2/0.75/14,7/4.0/0.75/15} {
  \node[gpu] (h\i) at (\x,\y) {\r};
}

\draw[fast] (g0) -- (g1);
\draw[fast] (g1) -- (g2);
\draw[fast] (g2) -- (g3);
\draw[fast] (g4) -- (g5);
\draw[fast] (g5) -- (g6);
\draw[fast] (g6) -- (g7);

\draw[fast] (h0) -- (h1);
\draw[fast] (h1) -- (h2);
\draw[fast] (h2) -- (h3);
\draw[fast] (h4) -- (h5);
\draw[fast] (h5) -- (h6);
\draw[fast] (h6) -- (h7);

\draw[slow] (g3) -- node[above,font=\scriptsize] {cross-node} (h0);

\node[align=left, text width=10.8cm] at (0,-1.2) {
  TP 通信通常最高频，应尽量限制在节点内高速互联范围。
  如果 TP group 跨节点，每层都会付出跨节点通信成本，扩展效率可能急剧下降。
};
```

\end{tikzpicture}
\caption{rank mapping 对通信路径的影响}
\label{fig:rank-mapping-topology}
\end{figure}

\subsection{为什么 rank mapping 影响性能}
\label{subsec:why-rank-mapping-affects-performance}

通信性能不是只由总 GPU 数决定，而是由\textbf{哪个 rank 和哪个 rank 通信、通信走哪条链路、链路是否拥塞}决定。

TP 通信常发生在每层内部。一个 80 层模型，每层 attention output 和 FFN down projection 至少各一次 TP all-reduce，则一个 forward/backward step 中 TP collective 次数非常多。如果 TP group 跨节点，每层都要跨网络同步，通信尾部可能无法隐藏。

PP 通信发生在 stage 边界。若 micro-batch 数 $M$ 很大，stage 间 send/recv 次数为 $O(MP_{\mathrm{PP}})$。如果相邻 PP stages 被放在跨机架或拥塞链路上，pipeline bubble 和等待时间会上升。

DP 通信通常发生在 backward 梯度同步。DDP bucket 可以与 backward 重叠，ZeRO/FSDP 也能通过 reduce-scatter/all-gather 改变通信形态。DP 跨节点是常见选择，但仍需要高带宽网络。

因此，一个经验性 rank mapping 原则是：

[
\text{通信频率越高、延迟越敏感的 group，越应该放在更近的拓扑范围内。}
]

按优先级粗略排序：

[
\text{TP}

>

\text{PP adjacent stages}

>

\text{DP}.
]

但这不是绝对。若 PP activation 很大、micro-batch 很多，PP 也可能成为主要瓶颈。若 ZeRO-3/FSDP 频繁 all-gather 参数，DP 维通信也会非常重。真正配置时需要结合 profiler、NCCL tests 和训练 step breakdown。

\section{计算、显存与通信的三角权衡}
\label{sec:compute-memory-communication-triangle}

3D 并行配置的本质，是在计算、显存与通信之间做权衡。一个配置可能显存很省，但通信太重；另一个配置计算效率高，但单卡显存放不下；还有一个配置通信少，但 pipeline bubble 大。

我们可以把三者画成一个三角形：

[
\text{Compute}
\leftrightarrow
\text{Memory}
\leftrightarrow
\text{Communication}.
]

\begin{figure}[H]
\centering
\begin{tikzpicture}[
font=\small,
vertex/.style={circle, draw, minimum size=1.25cm, align=center},
arrow/.style={-{Latex[length=2.2mm]}, thick}
]

```
\node[font=\bfseries] at (0,3.2) {3D 并行的 Compute-Memory-Communication 三角权衡};

\node[vertex, fill=blue!8] (compute) at (0,1.7) {Compute\\效率};
\node[vertex, fill=green!10] (memory) at (-4.0,-1.2) {Memory\\显存};
\node[vertex, fill=orange!14] (comm) at (4.0,-1.2) {Communication\\通信};

\draw[arrow] (memory) -- node[left, align=center] {切分模型\\降低显存} (compute);
\draw[arrow] (compute) -- node[right, align=center] {GEMM 变小\\通信增加} (comm);
\draw[arrow] (comm) -- node[below, align=center] {通信 buffer\\拓扑限制} (memory);

\node[align=left, text width=10.5cm] at (0,-2.7) {
  DP、TP、PP 的配置不是单目标优化。TP 增大可降低层内显存但增加高频通信；
  PP 增大可降低层数相关显存但增加 bubble 和 activation send/recv；
  DP 增大可提升数据吞吐但增加梯度同步压力。
};
```

\end{tikzpicture}
\caption{计算、显存与通信的三角权衡}
\label{fig:compute-memory-communication-triangle}
\end{figure}

\subsection{显存占用估算}
\label{subsec:memory-estimation}

训练显存可以拆成五部分：

[
M_{\mathrm{total}}
==================

M_{\mathrm{params}}
+
M_{\mathrm{grads}}
+
M_{\mathrm{optimizer}}
+
M_{\mathrm{activations}}
+
M_{\mathrm{buffers}}.
]

对于一个 dense Transformer，如果暂时忽略 embedding/lm head 和不均匀层切分，模型层参数在 TP 和 PP 下的每卡占用近似为：

[
M_{\mathrm{params,perGPU}}
\approx
\frac{N_{\mathrm{params}}\cdot s_{\mathrm{param}}}
{P_{\mathrm{TP}}P_{\mathrm{PP}}}.
]

这里的解释是：TP 切分层内矩阵，PP 切分层数。DP 不减少每个副本中的参数占用，除非使用 ZeRO/FSDP。

梯度与参数 shard 对应，若不使用 ZeRO/FSDP：

[
M_{\mathrm{grads,perGPU}}
\approx
\frac{N_{\mathrm{params}}\cdot s_{\mathrm{grad}}}
{P_{\mathrm{TP}}P_{\mathrm{PP}}}.
]

AdamW 优化器状态通常包括 master weight、一阶矩和二阶矩。如果它们随 TP/PP shard 保存，则：

[
M_{\mathrm{optimizer,perGPU}}
\approx
\frac{N_{\mathrm{params}}\cdot s_{\mathrm{optim}}}
{P_{\mathrm{TP}}P_{\mathrm{PP}}}.
]

如果使用 ZeRO-1 在 DP 维切分 optimizer state，则优化器状态进一步除以 $P_{\mathrm{DP}}$：

[
M_{\mathrm{optimizer,perGPU}}^{\mathrm{ZeRO1}}
\approx
\frac{N_{\mathrm{params}}\cdot s_{\mathrm{optim}}}
{P_{\mathrm{TP}}P_{\mathrm{PP}}P_{\mathrm{DP}}}.
]

如果使用 ZeRO-2，梯度也可在 DP 维切分；ZeRO-3 则参数也可在 DP 维切分。但 ZeRO-3 与 TP/PP 组合后通信更加复杂，后续 ZeRO 章节会专门讨论。

激活显存的估算更复杂。对一个 Transformer 层，隐藏状态激活基础大小为：

[
B_{\mathrm{micro}}S H s_{\mathrm{act}}.
]

若每个 PP stage 有 $L/P_{\mathrm{PP}}$ 层，则每个 stage 的激活显存粗略与：

[
M_{\mathrm{act,stage}}
\propto
N_{\mathrm{live\ microbatches}}
\cdot
\frac{L}{P_{\mathrm{PP}}}
\cdot
B_{\mathrm{micro}}S Hs_{\mathrm{act}}
]

相关。其中 $N_{\mathrm{live\ microbatches}}$ 与 pipeline schedule 有关。GPipe 中它可能接近 $M$；1F1B 中通常显著更低。

TP 对激活显存的影响取决于是否使用 sequence parallel。若仅使用 TP，一些中间激活沿 hidden 维切分，但 residual stream 可能仍 replicated。使用 sequence parallel 后，部分激活可进一步除以 $P_{\mathrm{TP}}$。

因此，真实估算可以写成：

[
M_{\mathrm{act}}
================

M_{\mathrm{act,sharded}}
+
M_{\mathrm{act,replicated}}
+
M_{\mathrm{act,temp}}.
]

如果粗略引入一个激活切分因子 $q_{\mathrm{act}}$：

[
M_{\mathrm{act,perGPU}}
\approx
\frac{
c_{\mathrm{act}}
\cdot
N_{\mathrm{live\ microbatches}}
\cdot
\frac{L}{P_{\mathrm{PP}}}
\cdot
B_{\mathrm{micro}}S Hs_{\mathrm{act}}
}
{q_{\mathrm{act}}},
]

其中 $c_{\mathrm{act}}$ 表示每层实际保存激活相对 hidden state 的倍数，$q_{\mathrm{act}}$ 可能为 $1$、$P_{\mathrm{TP}}$ 或介于两者之间。

下面给出一个简化显存估算代码，用于配置初筛。

\begin{lstlisting}[language=Python,caption={3D 并行训练显存的粗略估算器},label={lst:3d-memory-estimator}]
def estimate_3d_memory_gb(
num_params,
num_layers,
hidden_size,
seq_len,
micro_batch_size,
tp_size,
pp_size,
dp_size,
live_microbatches,
param_bytes=2,      # BF16 parameters
grad_bytes=2,       # BF16 gradients
master_bytes=4,     # FP32 master weights
adam_m_bytes=4,     # FP32 Adam m
adam_v_bytes=4,     # FP32 Adam v
act_bytes=2,        # BF16 activations
activation_factor=8.0,
activation_shard_factor=1.0,
zero_stage=0,
buffer_gb=4.0,
):
"""
A rough estimator for screening configurations.
It is not a replacement for real profiler-based measurement.
"""

```
# Parameters are sharded by TP and PP.
param_shard = num_params / (tp_size * pp_size)

param_mem = param_shard * param_bytes
grad_mem = param_shard * grad_bytes

optim_bytes = master_bytes + adam_m_bytes + adam_v_bytes
optim_mem = param_shard * optim_bytes

# ZeRO effects along DP dimension.
if zero_stage >= 1:
    optim_mem /= dp_size
if zero_stage >= 2:
    grad_mem /= dp_size
if zero_stage >= 3:
    param_mem /= dp_size

layers_per_stage = num_layers / pp_size
hidden_act = micro_batch_size * seq_len * hidden_size * act_bytes

act_mem = (
    live_microbatches
    * layers_per_stage
    * hidden_act
    * activation_factor
    / activation_shard_factor
)

total = param_mem + grad_mem + optim_mem + act_mem + buffer_gb * (1024 ** 3)

return {
    "params_gb": param_mem / (1024 ** 3),
    "grads_gb": grad_mem / (1024 ** 3),
    "optimizer_gb": optim_mem / (1024 ** 3),
    "activations_gb": act_mem / (1024 ** 3),
    "buffers_gb": buffer_gb,
    "total_gb": total / (1024 ** 3),
}
```

memory = estimate_3d_memory_gb(
num_params=70e9,
num_layers=80,
hidden_size=8192,
seq_len=4096,
micro_batch_size=1,
tp_size=8,
pp_size=4,
dp_size=8,
live_microbatches=4,
activation_shard_factor=8,  # assume sequence parallel roughly helps
zero_stage=1,
)

for k, v in memory.items():
print(k, round(v, 2))
\end{lstlisting}

这段代码只能用于直觉估算。真实显存还会受到 allocator fragmentation、kernel workspace、CUDA graph、NCCL buffer、FlashAttention workspace、recompute 策略、模型具体结构、MoE、embedding/lm head、vocab parallel cross entropy 等影响。

\subsection{通信量估算}
\label{subsec:communication-estimation}

3D 并行中的通信可以分为三类：

[
C_{\mathrm{comm}}
=================

C_{\mathrm{DP}}
+
C_{\mathrm{TP}}
+
C_{\mathrm{PP}}.
]

\textbf{DP 通信}主要是梯度同步或 ZeRO/FSDP 的状态同步。如果不使用 ZeRO，DDP 每 step 需要同步参数梯度。假设每个 DP replica 内部经过 TP/PP 后，每个 rank 持有的参数 shard 大小为：

[
G_{\mathrm{rank}}
=================

\frac{N_{\mathrm{params}}s_{\mathrm{grad}}}
{P_{\mathrm{TP}}P_{\mathrm{PP}}}.
]

普通 all-reduce 的每 rank 通信量可粗略按 Ring AllReduce 估算：

[
V_{\mathrm{DP}}
\approx
\frac{2(P_{\mathrm{DP}}-1)}{P_{\mathrm{DP}}}
G_{\mathrm{rank}}.
]

若使用 ZeRO-2 或 ZeRO-3，通信形态会变成 reduce-scatter、all-gather 和参数预取，不能简单用 DDP all-reduce 代替。

\textbf{TP 通信}发生在每层内部。Megatron 风格每层至少有 attention output projection 后一次 all-reduce、FFN down projection 后一次 all-reduce。每次 all-reduce 的张量大小大致为：

[
B_{\mathrm{micro}}S Hs_{\mathrm{act}}.
]

如果每个 stage 有 $L/P_{\mathrm{PP}}$ 层，每个 micro-batch 在本 stage 触发约 $n_{\mathrm{tp_collectives_per_layer}}$ 次 TP collective，则每个 step TP 通信量与：

[
M
\cdot
\frac{L}{P_{\mathrm{PP}}}
\cdot
n_{\mathrm{tp}}
\cdot
B_{\mathrm{micro}}S Hs_{\mathrm{act}}
]

相关。再乘以 all-reduce 的通信放大系数：

[
\frac{2(P_{\mathrm{TP}}-1)}{P_{\mathrm{TP}}}.
]

因此：

[
V_{\mathrm{TP}}
\approx
M
\cdot
\frac{L}{P_{\mathrm{PP}}}
\cdot
n_{\mathrm{tp}}
\cdot
\frac{2(P_{\mathrm{TP}}-1)}{P_{\mathrm{TP}}}
\cdot
B_{\mathrm{micro}}S Hs_{\mathrm{act}}.
]

\textbf{PP 通信}是 stage 边界传递 activation 和 activation gradient。每条边界每个 micro-batch forward 发送一次激活，backward 发送一次梯度：

[
V_{\mathrm{PP,boundary}}
\approx
2M B_{\mathrm{micro}}S Hs_{\mathrm{act}}.
]

每个 rank 只和相邻 stage 通信，不一定经过所有边界。因此每 rank PP 通信取决于它是否在中间 stage、首 stage 还是尾 stage。中间 stage 既接收又发送，边界通信更多；首尾 stage 少一侧。

通信估算的重点不是得到一个完美数字，而是判断哪个维度最可能成为瓶颈：

[
\text{TP 高频，偏低延迟和高带宽节点内通信；}
]
[
\text{PP 与 micro-batch 和激活大小相关，偏点对点通信；}
]
[
\text{DP 与参数/梯度大小相关，偏大块 collective 通信。}
]

\begin{table}[H]
\centering
\small
\caption{DP、TP、PP 通信模式对比}
\label{tab:dp-tp-pp-communication-comparison}
\begin{tabular}{p{0.16\textwidth}p{0.25\textwidth}p{0.24\textwidth}p{0.25\textwidth}}
\toprule
\textbf{并行维度} & \textbf{通信对象} & \textbf{常见操作} & \textbf{性能敏感点} \
\midrule
DP
& 梯度、参数 shard、优化器状态
& all-reduce, reduce-scatter, all-gather
& 跨节点带宽、bucket overlap、ZeRO stage。 \
\midrule
TP
& 层内激活、partial output、sequence shard
& all-reduce, all-gather, reduce-scatter
& 频率高，强依赖节点内 NVLink/NVSwitch。 \
\midrule
PP
& stage 边界激活与激活梯度
& send/recv
& micro-batch 数、hidden/seq 大小、stage 拓扑。 \
\midrule
MoE EP
& token dispatch 和 expert output
& all-to-all
& 负载均衡、网络拥塞、expert placement。 \
\bottomrule
\end{tabular}
\end{table}

\subsection{计算效率估算}
\label{subsec:compute-efficiency-estimation}

大模型训练常用 MFU，Model FLOPs Utilization，衡量有效模型计算占硬件峰值计算的比例：

[
\mathrm{MFU}
============

\frac{\mathrm{Model\ FLOPs\ per\ step}}
{\mathrm{Step\ time}\cdot \mathrm{Peak\ FLOPs\ of\ GPUs}}.
]

如果模型每个 token 的训练 FLOPs 粗略为 $6N$，每 step 训练 token 数为：

[
T_{\mathrm{tokens/step}}
========================

B_{\mathrm{global}}S,
]

则每 step 模型 FLOPs 为：

[
F_{\mathrm{step}}
\approx
6N B_{\mathrm{global}}S.
]

若使用 $W$ 张 GPU，每张 GPU 峰值为 $F_{\mathrm{peak}}$，step time 为 $T_{\mathrm{step}}$，则：

[
\mathrm{MFU}
\approx
\frac{6N B_{\mathrm{global}}S}
{T_{\mathrm{step}} W F_{\mathrm{peak}}}.
]

MFU 下降可能来自：

\begin{itemize}
\item TP size 太大导致 GEMM 太小；
\item PP bubble 占比高；
\item DP 通信暴露；
\item 数据加载阻塞；
\item checkpoint 阻塞；
\item kernel fusion 不足，小 kernel 太多；
\item activation checkpointing 重计算太多；
\item straggler 或节点性能不一致；
\item 通信与计算资源争抢。
\end{itemize}

Pipeline bubble 对计算效率的影响可以单独估算。若使用 $M$ 个 micro-batches、$P_{\mathrm{PP}}$ 个 stages，简单近似：

[
\eta_{\mathrm{pipe}}
\approx
\frac{M}{M+P_{\mathrm{PP}}-1}.
]

例如 $P_{\mathrm{PP}}=8$，$M=8$：

[
\eta_{\mathrm{pipe}}
====================

# \frac{8}{8+8-1}

\frac{8}{15}
\approx 53.3%.
]

这说明 micro-batch 数与 PP size 相当时，bubble 仍然很大。若 $M=64$：

[
\eta_{\mathrm{pipe}}
====================

# \frac{64}{64+8-1}

\frac{64}{71}
\approx 90.1%.
]

但 $M$ 变大意味着 gradient accumulation 变大，global batch size 也变大，可能影响训练动力学。因此 pipeline efficiency 不能孤立优化。

\begin{lstlisting}[language=Python,caption={MFU 与 Pipeline Efficiency 的简化估算},label={lst:mfu-pipeline-efficiency-estimator}]
def estimate_mfu(
num_params,
global_batch_size,
seq_len,
step_time_sec,
num_gpus,
peak_flops_per_gpu,
):
# Dense Transformer training FLOPs rough estimate.
flops_per_step = 6 * num_params * global_batch_size * seq_len
peak_flops_per_step = step_time_sec * num_gpus * peak_flops_per_gpu
return flops_per_step / peak_flops_per_step

def pipeline_efficiency(num_microbatches, pp_size):
return num_microbatches / (num_microbatches + pp_size - 1)

for m in [4, 8, 16, 32, 64]:
print(m, round(pipeline_efficiency(m, pp_size=8), 3))
\end{lstlisting}

\subsection{不同规模模型的并行策略选择}
\label{subsec:parallel-strategy-for-different-model-sizes}

并行策略没有唯一最优答案。下面给出一组工程直觉，用于初始配置选择。

\textbf{小模型，例如 1B 到 7B。}如果单卡可以放下模型和训练状态，优先使用 DDP 或 FSDP。TP/PP 可能降低单卡计算效率，除非 batch 或序列足够大。常见策略：

[
P_{\mathrm{TP}}=1,\quad
P_{\mathrm{PP}}=1,\quad
P_{\mathrm{DP}}=W.
]

如果显存不足，可以使用 ZeRO/FSDP、activation checkpointing 或 LoRA/QLoRA 等方法。

\textbf{中等模型，例如 7B 到 30B。}可能需要 TP 或 ZeRO。若节点内 8 卡互联较好，可尝试：

[
P_{\mathrm{TP}}=2\ \text{or}\ 4,
\quad
P_{\mathrm{PP}}=1,
\quad
P_{\mathrm{DP}}=\frac{W}{P_{\mathrm{TP}}}.
]

如果层数较多、显存仍不足，可引入 PP：

[
P_{\mathrm{TP}}=4,
\quad
P_{\mathrm{PP}}=2,
\quad
P_{\mathrm{DP}}=\frac{W}{8}.
]

\textbf{大模型，例如 30B 到 100B。}通常需要 TP+PP+DP。TP 放在节点内，PP 跨节点或跨节点组，DP 复制多个 3D replica。典型方向：

[
P_{\mathrm{TP}}=4\ \text{or}\ 8,
\quad
P_{\mathrm{PP}}=4\ \text{or}\ 8,
\quad
P_{\mathrm{DP}}=\frac{W}{P_{\mathrm{TP}}P_{\mathrm{PP}}}.
]

\textbf{超大模型，例如 100B 以上。}通常需要 3D 并行加 ZeRO、activation checkpointing、sequence parallel、distributed checkpoint 和拓扑感知调度。TP 可能取节点内最大高速互联规模，例如 8；PP 取决于层数和显存；DP size 由总 GPU 数和 global batch 需求决定。

\begin{table}[H]
\centering
\small
\caption{不同模型规模下的并行策略直觉}
\label{tab:parallel-strategy-by-model-size}
\begin{tabular}{p{0.18\textwidth}p{0.25\textwidth}p{0.24\textwidth}p{0.23\textwidth}}
\toprule
\textbf{模型规模} & \textbf{优先策略} & \textbf{常见配置方向} & \textbf{主要风险} \
\midrule
1B--7B
& DDP / FSDP
& TP=1, PP=1，扩大 DP
& 显存、数据加载、小 batch 通信。 \
\midrule
7B--30B
& TP 或 ZeRO/FSDP
& TP=2/4，PP 可为 1/2
& TP 通信、GEMM 变小、ZeRO 通信。 \
\midrule
30B--100B
& TP+PP+DP
& TP=4/8，PP=2/4/8
& pipeline bubble、stage 均衡、rank mapping。 \
\midrule
100B+
& 3D + ZeRO + checkpointing
& TP 节点内，PP 跨节点，DP 多副本
& 通信、checkpoint、故障恢复、长时间稳定性。 \
\bottomrule
\end{tabular}
\end{table}

这些只是初始猜测。真实项目必须通过 profile 和试跑验证。一个配置是否好，不看理论是否优雅，而看：

\begin{itemize}
\item 单卡显存是否有安全余量；
\item MFU 是否合理；
\item step time 是否稳定；
\item pipeline bubble 是否可接受；
\item NCCL 时间占比是否过高；
\item checkpoint 是否能在可接受时间内完成；
\item 故障恢复是否可靠；
\item global batch 和学习率是否满足训练 recipe。
\end{itemize}

\section{大模型训练配置推导}
\label{sec:large-model-training-config-derivation}

3D 并行配置推导是 AI Infra 工程师的核心能力之一。给定模型规模、GPU 数量、单卡显存、节点拓扑和目标 global batch，需要反推出 TP、PP、DP、micro-batch、gradient accumulation 等参数。

\subsection{global batch size}
\label{subsec:global-batch-size}

在 3D 并行中，global batch size 由数据并行、副本内 micro-batch 和梯度累积共同决定：

[
B_{\mathrm{global}}
===================

B_{\mathrm{micro}}
\cdot
M
\cdot
P_{\mathrm{DP}}.
]

其中：

\begin{itemize}
\item $B_{\mathrm{micro}}$ 是每个 pipeline micro-batch 的 batch size；
\item $M$ 是每个 optimizer step 的 micro-batch 数；
\item $P_{\mathrm{DP}}$ 是数据并行副本数。
\end{itemize}

如果按 token 数计算：

[
T_{\mathrm{global}}
===================

B_{\mathrm{global}}
\cdot
S
=

B_{\mathrm{micro}}
\cdot
M
\cdot
P_{\mathrm{DP}}
\cdot
S.
]

在语言模型训练中，常用 tokens per step 描述 batch，因为序列长度直接影响计算量。若使用 sequence packing，实际每个 sample 的有效 token 可能不同，此时需要按 non-padding token 数统计。

\subsection{micro batch size}
\label{subsec:micro-batch-size}

Micro-batch size 是 3D 并行中非常敏感的参数。它影响：

\begin{itemize}
\item 单个 micro-batch 激活显存；
\item GEMM shape 和 Tensor Core 利用率；
\item pipeline bubble 摊薄方式；
\item gradient accumulation steps；
\item global batch size；
\item data loader 和 sequence packing。
\end{itemize}

通常显存越紧，$B_{\mathrm{micro}}$ 越小。很多大模型训练中，micro-batch size 甚至可能为 $1$。这并不意味着 global batch size 小，因为可以通过更大的 $M$ 和 $P_{\mathrm{DP}}$ 累积到目标 global batch。

但是，$B_{\mathrm{micro}}=1$ 也可能带来问题：

\begin{itemize}
\item GEMM 的 $M$ 维，即 token 数 $B_{\mathrm{micro}}S$，可能较小；
\item pipeline micro-batch 数需要增加，调度开销上升；
\item 如果序列长度也不大，GPU 利用率会下降；
\item Batch 维相关的数据增强或采样统计变弱。
\end{itemize}

因此配置时通常先从显存允许的最大 micro-batch 试起，再根据 global batch 需求调 gradient accumulation。

\subsection{gradient accumulation}
\label{subsec:gradient-accumulation-in-3d}

在 PP 中，micro-batch 数 $M$ 既是 gradient accumulation 的一部分，也是 pipeline 填充效率的关键。为了达到目标 global batch：

[
M
=

\frac{B_{\mathrm{global}}}
{B_{\mathrm{micro}}P_{\mathrm{DP}}}.
]

同时为了降低 pipeline bubble，希望：

[
M \gg P_{\mathrm{PP}}.
]

这两个条件可能冲突。如果 $P_{\mathrm{DP}}$ 很大，给定 global batch 下 $M$ 可能变小，导致 pipeline bubble 大。反过来，如果为了提高 pipeline efficiency 增大 $M$，global batch 可能过大，需要调整学习率或训练 recipe。

例如：

[
B_{\mathrm{global}}=1024,
\quad
B_{\mathrm{micro}}=1,
\quad
P_{\mathrm{DP}}=64.
]

则：

[
M=\frac{1024}{1\times 64}=16.
]

若：

[
P_{\mathrm{PP}}=16,
]

则 pipeline efficiency 近似为：

[
\eta_{\mathrm{pipe}}
====================

# \frac{16}{16+16-1}

\frac{16}{31}
\approx 51.6%.
]

这很低。可以考虑：

\begin{itemize}
\item 减小 DP size，增大 $M$；
\item 减小 PP size；
\item 增大 global batch；
\item 增大 micro-batch；
\item 使用 interleaved pipeline；
\item 调整并行切分，使 PP bubble 降低。
\end{itemize}

\subsection{TP size、PP size、DP size}
\label{subsec:tp-pp-dp-size}

选择 TP、PP、DP size 时，可以按以下顺序思考。

\textbf{第一步：选择 TP size。}

TP size 主要受以下因素约束：

\begin{itemize}
\item hidden size 是否能整除 TP size；
\item attention heads 是否能整除 TP size；
\item GQA 中 kv heads 是否能合理切分；
\item TP group 是否能放进节点内高速互联；
\item 单层矩阵是否太大，需要更多 TP；
\item TP 后 GEMM 是否仍然足够大。
\end{itemize}

常见策略是 TP size 选择节点内 GPU 数的因子，例如 2、4、8。对于 8 卡 NVSwitch 节点，TP=8 是常见上限。若 TP 跨节点，需要非常谨慎。

\textbf{第二步：选择 PP size。}

PP size 主要受以下因素约束：

\begin{itemize}
\item 模型层数 $L$ 是否能较均匀切分；
\item 每个 stage 的参数和激活显存是否放得下；
\item PP bubble 是否可接受；
\item micro-batch 数是否足够大；
\item stage 间通信是否会跨慢链路；
\item embedding/lm head 是否造成首尾 stage 不均。
\end{itemize}

一般来说，PP size 越大，每个 stage 层数越少，显存越省；但 bubble 和调度复杂度越高。PP size 也不宜超过层数的合理切分粒度。

\textbf{第三步：由总 GPU 数反推 DP size。}

当 $P_{\mathrm{TP}}$ 和 $P_{\mathrm{PP}}$ 确定后：

[
P_{\mathrm{DP}}
===============

\frac{W}{P_{\mathrm{TP}}P_{\mathrm{PP}}}.
]

DP size 决定可并行处理多少数据副本，也影响 global batch 和梯度同步规模。若 DP size 太小，global batch 可能不足或训练吞吐不足；若 DP size 太大，在固定 global batch 下 micro-batch 数可能太小，pipeline bubble 上升。

\subsection{从 GPU 数量反推并行配置}
\label{subsec:derive-parallel-config-from-gpu-count}

下面给出一个配置推导流程。假设：

\begin{itemize}
\item 总 GPU 数 $W=256$；
\item 每节点 8 张 GPU；
\item 模型有 $L=80$ 层；
\item hidden size $H=8192$；
\item attention heads $n_h=64$；
\item 目标 global batch size $B_{\mathrm{global}}=1024$；
\item 序列长度 $S=4096$；
\item 初始设 micro-batch size $B_{\mathrm{micro}}=1$。
\end{itemize}

\textbf{候选配置 A：}

[
P_{\mathrm{TP}}=8,\quad
P_{\mathrm{PP}}=4.
]

则：

[
P_{\mathrm{DP}}
===============

# \frac{256}{8\times 4}

8.

]

Gradient accumulation micro-batch 数为：

[
M
=

# \frac{1024}{1\times 8}

128.

]

Pipeline efficiency 近似为：

[
\eta_{\mathrm{pipe}}
====================

# \frac{128}{128+4-1}

\frac{128}{131}
\approx 97.7%.
]

这个配置 pipeline bubble 很低。TP=8 可放在每个节点内；PP=4 意味着一个 DP replica 需要：

[
8\times 4=32
]

张 GPU，也就是 4 个节点。总共 256 GPU 可形成 8 个 DP replicas。

\textbf{候选配置 B：}

[
P_{\mathrm{TP}}=4,\quad
P_{\mathrm{PP}}=8.
]

则：

[
P_{\mathrm{DP}}
===============

# \frac{256}{4\times 8}

8.

]

Micro-batch 数仍为：

[
M=128.
]

Pipeline efficiency：

[
\eta_{\mathrm{pipe}}
====================

# \frac{128}{128+8-1}

\frac{128}{135}
\approx 94.8%.
]

这个配置 TP 通信更少，单个 TP group 只用 4 张 GPU；PP 更深，每个 stage 层数约 $10$。如果显存允许，这可能比 TP=8 更好，因为 GEMM 更大、TP 通信更少。但 PP stage 间通信和调度稍多。

\textbf{候选配置 C：}

[
P_{\mathrm{TP}}=8,\quad
P_{\mathrm{PP}}=8.
]

则：

[
P_{\mathrm{DP}}
===============

# \frac{256}{8\times 8}

4.

]

Micro-batch 数：

[
M
=

# \frac{1024}{1\times 4}

256.

]

Pipeline efficiency：

[
\eta_{\mathrm{pipe}}
====================

# \frac{256}{256+8-1}

\frac{256}{263}
\approx 97.3%.
]

这个配置显存更省，但每个 DP replica 需要 64 GPU，只有 4 个 DP replicas。它适合模型更大、显存更紧的情况，但 TP 和 PP 都较高，通信和调度复杂度更大。

下面给出一个枚举候选配置的脚本。

\begin{lstlisting}[language=Python,caption={从 GPU 数量枚举 3D 并行候选配置},label={lst:enumerate-3d-configs}]
def divisors(n):
return [x for x in range(1, n + 1) if n % x == 0]

def enumerate_3d_configs(
world_size,
gpus_per_node,
num_layers,
num_heads,
hidden_size,
global_batch_size,
micro_batch_size,
preferred_tp_values=(1, 2, 4, 8),
):
configs = []

```
for tp in preferred_tp_values:
    if world_size % tp != 0:
        continue
    if num_heads % tp != 0:
        continue
    if hidden_size % tp != 0:
        continue

    # Prefer TP within a node.
    if tp > gpus_per_node:
        continue

    remaining = world_size // tp

    for pp in divisors(remaining):
        if num_layers % pp != 0:
            continue

        dp = remaining // pp
        if dp <= 0:
            continue

        denom = micro_batch_size * dp
        if global_batch_size % denom != 0:
            continue

        microbatches = global_batch_size // denom
        pipe_eff = microbatches / (microbatches + pp - 1)

        configs.append({
            "tp": tp,
            "pp": pp,
            "dp": dp,
            "microbatches": microbatches,
            "layers_per_stage": num_layers // pp,
            "pipeline_efficiency": pipe_eff,
            "gpus_per_replica": tp * pp,
        })

# Sort by high pipeline efficiency, then smaller TP/PP complexity.
configs.sort(
    key=lambda c: (
        -c["pipeline_efficiency"],
        c["tp"] * c["pp"],
        c["pp"],
    )
)
return configs
```

configs = enumerate_3d_configs(
world_size=256,
gpus_per_node=8,
num_layers=80,
num_heads=64,
hidden_size=8192,
global_batch_size=1024,
micro_batch_size=1,
)

for cfg in configs[:10]:
print(cfg)
\end{lstlisting}

这个脚本只做基本整除和 pipeline efficiency 检查。真实系统还需要把显存估算、通信估算、节点拓扑、ZeRO stage、checkpoint 策略和训练 recipe 一起纳入搜索。

\section{故障恢复与 checkpoint}
\label{sec:fault-tolerance-and-checkpoint}

大模型训练通常持续数天、数周甚至更久。随着 GPU 数量增加，故障概率也会增加。假设单个节点在训练期间发生故障的概率为 $p$，节点数为 $N_{\mathrm{nodes}}$，整个作业无故障概率近似为：

[
(1-p)^{N_{\mathrm{nodes}}}.
]

当节点数很大时，即使单节点故障概率很低，整个作业也很可能遇到故障。因此，大规模训练系统必须具备 checkpoint 和恢复能力。

Checkpoint 不只是保存模型权重。一个完整训练 checkpoint 通常包括：

\begin{itemize}
\item 模型参数；
\item 优化器状态；
\item 梯度缩放器状态；
\item 学习率 scheduler 状态；
\item global step；
\item consumed samples 或 consumed tokens；
\item tokenizer 和数据状态；
\item RNG states；
\item 并行配置和 rank mapping；
\item ZeRO/FSDP/TP/PP shard metadata；
\item dataloader 或 dataset cursor；
\item 训练配置文件与代码版本信息。
\end{itemize}

如果少保存其中某些状态，训练可能“能恢复”，但不是严格恢复。比如 scheduler step 错了，学习率会跳；consumed tokens 错了，数据会重复或跳过；RNG state 错了，dropout 和数据 shuffle 不一致；rank mapping metadata 错了，sharded 权重可能加载到错误 rank。

\subsection{distributed checkpoint}
\label{subsec:distributed-checkpoint}

单机小模型常由 rank 0 保存完整 checkpoint。但 3D 并行训练中，这通常不可行。原因包括：

\begin{itemize}
\item 单个 rank 不一定拥有完整模型；
\item TP/PP/ZeRO/FSDP 后参数和优化器状态是分片的；
\item checkpoint 体积可能达到 TB 级；
\item rank 0 聚合完整状态会造成显存和内存峰值；
\item 单点写入会造成 IO 瓶颈；
\item 恢复时重新广播完整状态代价巨大。
\end{itemize}

Distributed checkpoint 的思想是：每个 rank 或每组 rank 保存自己负责的 shard，形成一个分布式 checkpoint 目录。目录中包含：

\begin{itemize}
\item 每个 shard 的数据文件；
\item 描述全局张量如何切分的 metadata；
\item rank 到 shard 的映射；
\item tensor 名称、shape、dtype、offset；
\item 训练状态元数据。
\end{itemize}

恢复时，每个 rank 根据当前并行配置和 metadata 读取自己需要的 shard。如果恢复配置与保存配置相同，读取较简单；如果 TP/PP/DP size 改变，则需要做 resharding。

\begin{figure}[H]
\centering
\begin{tikzpicture}[
font=\small,
rank/.style={draw, rounded corners=3pt, minimum width=1.45cm, minimum height=0.65cm, align=center, fill=blue!6},
file/.style={draw, rounded corners=2pt, minimum width=1.35cm, minimum height=0.5cm, align=center, fill=orange!14},
meta/.style={draw, rounded corners=3pt, minimum width=2.2cm, minimum height=0.65cm, align=center, fill=green!10},
arrow/.style={-{Latex[length=2.1mm]}, thick}
]

```
\node[font=\bfseries] at (0,3.35) {分布式 Checkpoint：每个 rank 保存自己的 shard};

\foreach \i/\x in {0/-4.8,1/-2.4,2/0,3/2.4,4/4.8} {
  \node[rank] (r\i) at (\x,1.55) {Rank \i};
  \node[file] (f\i) at (\x,0.45) {shard \i};
  \draw[arrow] (r\i) -- (f\i);
}

\node[meta] (meta) at (0,-0.8) {metadata.json\\tensor layout};
\foreach \i in {0,1,2,3,4} {
  \draw[arrow] (f\i) -- (meta);
}

\node[align=left, text width=10.6cm] at (0,-2.25) {
  分布式 checkpoint 不聚合到单个 rank，而是保存 shard 与 metadata。
  metadata 描述每个全局张量如何被 TP、PP、DP、ZeRO 等维度切分。
};
```

\end{tikzpicture}
\caption{分布式 checkpoint 保存流程}
\label{fig:distributed-checkpoint-flow}
\end{figure}

\subsection{sharded checkpoint}
\label{subsec:sharded-checkpoint}

Sharded checkpoint 是分布式训练中的常见格式。以一个线性层权重为例：

[
\mat{W}\in\mathbb{R}^{H_{\mathrm{in}}\times H_{\mathrm{out}}}.
]

如果它在 TP 维做 column parallel，rank $t$ 保存：

[
\mat{W}_t
=========

\mat{W}*{:,\ t\frac{H*{\mathrm{out}}}{P_{\mathrm{TP}}}:(t+1)\frac{H_{\mathrm{out}}}{P_{\mathrm{TP}}}}.
]

Checkpoint metadata 需要记录：

\begin{itemize}
\item tensor name；
\item global shape；
\item shard shape；
\item shard offset；
\item dtype；
\item parallel dimension；
\item owning rank；
\item file path。
\end{itemize}

一个简化 metadata 可以写成：

\begin{lstlisting}[language=Python,caption={Sharded checkpoint metadata 的简化示意},label={lst:sharded-checkpoint-metadata}]
metadata = {
"global_step": 120000,
"consumed_tokens": 503316480000,
"parallel_config": {
"tp_size": 8,
"pp_size": 4,
"dp_size": 8,
"zero_stage": 1,
},
"tensors": {
"layers.0.attention.qkv.weight": [
{
"rank": 0,
"global_shape": [8192, 24576],
"local_shape": [8192, 3072],
"offset": [0, 0],
"dtype": "bfloat16",
"file": "rank_00000_model.pt",
},
{
"rank": 1,
"global_shape": [8192, 24576],
"local_shape": [8192, 3072],
"offset": [0, 3072],
"dtype": "bfloat16",
"file": "rank_00001_model.pt",
},
]
}
}
\end{lstlisting}

如果后续想把 TP size 从 8 改为 4，需要把原来的 8 个 column shards 重新拼接再切成 4 个 shards。这就是 checkpoint resharding。它要求 metadata 足够完整，否则无法安全转换。

Sharded checkpoint 的优点是：

\begin{itemize}
\item 保存时无需聚合完整模型；
\item 每个 rank 并行写入，IO 吞吐更高；
\item 与 ZeRO/FSDP/TP/PP 的分片状态自然匹配；
\item 大模型 checkpoint 可扩展到 TB 级。
\end{itemize}

缺点是：

\begin{itemize}
\item 文件数量多；
\item metadata 更复杂；
\item resharding 困难；
\item 部分 shard 损坏会影响整体恢复；
\item 需要与对象存储或并行文件系统良好配合。
\end{itemize}

\subsection{optimizer state checkpoint}
\label{subsec:optimizer-state-checkpoint}

仅保存模型参数不足以恢复训练。AdamW 的优化器状态包括：

[
m_t,\quad v_t,\quad \theta_{\mathrm{master}}.
]

如果不保存这些状态，只用模型权重恢复，等价于重新初始化优化器。训练可能继续下降，但 loss 曲线、收敛路径和最终结果都会改变。

优化器状态 checkpoint 的挑战是体积大。对于 AdamW，优化器状态通常比参数本身大得多。若模型参数为 $N$，优化器状态可能为：

[
12N\ \mathrm{bytes}
]

甚至更多，取决于是否保存 master weight、dtype 和 ZeRO stage。对于 70B 模型，优化器状态可以达到数百 GB 甚至更高。

使用 ZeRO 或 FSDP 时，优化器状态通常在 DP 维切分。Checkpoint 时每个 DP rank 只保存自己负责的 optimizer shard。恢复时必须保证：

[
\text{optimizer shard}
\leftrightarrow
\text{parameter shard}
]

正确对应。如果参数 resharding 了，优化器状态也要同步 reshard，否则会出现状态错配。

下面给出一个简化的 checkpoint 保存流程。

\begin{lstlisting}[language=Python,caption={保存模型与优化器 shard 的简化伪代码},label={lst:save-sharded-checkpoint-pseudocode}]
def save_sharded_checkpoint(
model,
optimizer,
scheduler,
step,
consumed_tokens,
rank,
checkpoint_dir,
parallel_context,
):
# Each rank saves only local shards.
model_state = collect_local_model_shards(model)
optim_state = collect_local_optimizer_shards(optimizer)

```
rank_state = {
    "model": model_state,
    "optimizer": optim_state,
    "scheduler": scheduler.state_dict(),
    "step": step,
    "consumed_tokens": consumed_tokens,
    "rng_state": get_rng_state(),
    "parallel_context": parallel_context.local_metadata(),
}

path = f"{checkpoint_dir}/rank_{rank:05d}.pt"
torch.save(rank_state, path)

# Rank 0 writes global metadata after all ranks finish writing shards.
torch.distributed.barrier()

if rank == 0:
    metadata = build_global_checkpoint_metadata(
        model=model,
        optimizer=optimizer,
        step=step,
        consumed_tokens=consumed_tokens,
        parallel_context=parallel_context,
    )
    write_json(metadata, f"{checkpoint_dir}/metadata.json")

torch.distributed.barrier()
```

\end{lstlisting}

真实 checkpoint 系统还需要处理原子性。不能让训练在 metadata 写完但部分 shard 没写完时误认为 checkpoint 可用。常见做法是：

\begin{itemize}
\item 写入临时目录；
\item 所有 rank 完成后写 success marker；
\item 最后原子 rename 或更新 latest 指针；
\item 恢复时只读取带完整 marker 的 checkpoint。
\end{itemize}

\subsection{checkpoint 与训练吞吐的冲突}
\label{subsec:checkpoint-vs-throughput}

Checkpoint 会中断或拖慢训练。设每隔 $K$ steps 保存一次 checkpoint，每次 checkpoint 花费 $T_{\mathrm{ckpt}}$，训练 step time 为 $T_{\mathrm{step}}$，则摊销开销为：

[
\rho_{\mathrm{ckpt}}
====================

\frac{T_{\mathrm{ckpt}}}{KT_{\mathrm{step}}}.
]

如果每 1000 steps 保存一次，每步 2 秒，checkpoint 花 10 分钟：

[
\rho_{\mathrm{ckpt}}
====================

# \frac{600}{1000\times 2}

30%.
]

这非常高。为了降低影响，可以：

\begin{itemize}
\item 降低 checkpoint 频率；
\item 使用异步 checkpoint；
\item 使用分布式并行写入；
\item 只保存必要状态，减少临时 checkpoint 体积；
\item 周期性保存 full checkpoint，期间保存 lightweight checkpoint；
\item 使用高吞吐并行文件系统或对象存储；
\item 将 checkpoint 写入与训练计算 overlap。
\end{itemize}

但 checkpoint 频率不能无限降低。故障后需要回退到最近 checkpoint。如果 checkpoint 间隔太长，一次故障会损失大量 GPU 时间。设平均故障间隔为 MTBF，checkpoint 间隔为 $K T_{\mathrm{step}}$，每次故障平均损失约为半个 checkpoint 间隔：

[
T_{\mathrm{lost}}
\approx
\frac{1}{2}K T_{\mathrm{step}}.
]

因此 checkpoint 策略需要在吞吐开销和故障损失之间平衡。

\begin{figure}[H]
\centering
\begin{tikzpicture}[
font=\small,
axis/.style={-{Latex[length=2.1mm]}, thick},
train/.style={draw, rounded corners=2pt, minimum width=1.0cm, minimum height=0.45cm, align=center, fill=blue!12},
ckpt/.style={draw, rounded corners=2pt, minimum width=0.75cm, minimum height=0.45cm, align=center, fill=orange!16},
fail/.style={draw, star, star points=5, minimum size=0.8cm, fill=red!20}
]

```
\node[font=\bfseries] at (0,3.1) {Checkpoint 间隔与故障回退成本};

\draw[axis] (-5.5,1.2) -- (5.5,1.2) node[right] {time};

\foreach \i/\x in {0/-4.5,1/-3.3,2/-2.1,3/-0.9,4/0.3,5/1.5,6/2.7,7/3.9} {
  \node[train] at (\x,1.2) {train};
}

\foreach \x in {-5.1,-1.5,2.1} {
  \node[ckpt] at (\x,2.0) {ckpt};
  \draw[dashed] (\x,1.55) -- (\x,1.2);
}

\node[fail] at (4.2,1.95) {};
\node[align=center] at (4.2,2.65) {failure};

\draw[red!70, thick, <->] (2.1,0.45) -- (4.2,0.45);
\node[red!70, align=center] at (3.15,0.05) {lost work\\since last checkpoint};

\node[align=left, text width=10.5cm] at (0,-1.35) {
  checkpoint 太频繁会降低吞吐；太稀疏会在故障后损失更多训练时间。
  大规模训练必须根据故障率、checkpoint 时长和训练成本选择合适间隔。
};
```

\end{tikzpicture}
\caption{checkpoint 间隔与故障恢复损失}
\label{fig:checkpoint-interval-failure-cost}
\end{figure}

\section{3D 并行配置实战检查清单}
\label{sec:3d-parallel-config-checklist}

为了把本章内容落到工程实践，下面给出一个 3D 并行配置检查清单。

\subsection{模型结构检查}
\label{subsec:model-structure-checklist}

配置前需要确认：

\begin{itemize}
\item hidden size 是否能整除 TP size；
\item attention heads 是否能整除 TP size；
\item kv heads 是否能整除 TP size 或支持对应 GQA 切分；
\item FFN intermediate size 是否能整除 TP size；
\item vocabulary size 是否需要 padding 到 TP 友好的倍数；
\item layer 数是否能被 PP size 合理切分；
\item embedding/lm head 是否共享权重；
\item 是否有 MoE 层，expert 数是否能与 expert parallel 匹配；
\item 是否使用 sequence parallel；
\item 是否使用 activation checkpointing。
\end{itemize}

\subsection{硬件拓扑检查}
\label{subsec:hardware-topology-checklist}

需要确认：

\begin{itemize}
\item 每节点 GPU 数；
\item 节点内 NVLink/NVSwitch/PCIe 拓扑；
\item 每节点 NIC 数；
\item GPU 到 NIC 的亲和性；
\item 跨节点网络带宽；
\item 是否多 rail；
\item rank mapping 是否让 TP group 尽量节点内；
\item PP 相邻 stage 是否落在合理网络距离；
\item DP group 是否均匀分布；
\item NCCL 是否识别正确拓扑。
\end{itemize}

\subsection{训练配置检查}
\label{subsec:training-config-checklist}

需要确认：

\begin{itemize}
\item $W=P_{\mathrm{DP}}P_{\mathrm{TP}}P_{\mathrm{PP}}$；
\item $B_{\mathrm{global}}=B_{\mathrm{micro}}MP_{\mathrm{DP}}$；
\item $M$ 是否足够大以降低 pipeline bubble；
\item micro-batch size 是否能放入显存；
\item global batch 是否符合训练 recipe；
\item 学习率是否按 global batch 调整；
\item warmup steps 是否按 token 或 step 正确设置；
\item gradient clipping 是否在正确维度上执行；
\item mixed precision 与 loss scaling 是否配置正确；
\item activation checkpointing 粒度是否合理。
\end{itemize}

\subsection{运行时检查}
\label{subsec:runtime-checklist}

训练启动后，需要观察：

\begin{itemize}
\item 每张 GPU 显存峰值是否有余量；
\item 每个 rank step time 是否一致；
\item pipeline stage 是否负载均衡；
\item TP collective 是否占比过高；
\item DP gradient sync 是否被 backward 隐藏；
\item PP send/recv 是否出现等待；
\item dataloader 是否阻塞；
\item loss 是否稳定下降；
\item gradient norm 是否异常；
\item 是否出现 NaN/inf；
\item checkpoint 是否成功写入并可恢复；
\item tokens per second 和 MFU 是否达到预期。
\end{itemize}

\section{本章小结}
\label{sec:chapter10-summary}

本章把数据并行、张量并行和流水线并行组合起来，系统介绍了 3D 并行策略。

\begin{itemize}
\item \textbf{3D 并行} 将 DP、TP、PP 三个维度正交组合，总 GPU 数满足 $W=P_{\mathrm{DP}}P_{\mathrm{TP}}P_{\mathrm{PP}}$。
\item \textbf{DP 切分数据}，不同副本处理不同 batch 分片，并同步梯度或分片状态。
\item \textbf{TP 切分层内张量}，例如 attention heads、QKV projection、FFN 中间维度，通信频率最高，通常应限制在节点内高速互联范围。
\item \textbf{PP 切分模型层}，不同 stages 负责不同层段，通过 micro-batches 填充流水线，通信主要是相邻 stage 的 activation send/recv。
\item \textbf{group 划分} 是 3D 并行的实现基础。每个 rank 同时属于 TP group、PP group 和 DP group。
\item \textbf{rank mapping} 直接影响通信路径。高频通信的 group 应优先映射到更近、更快的物理拓扑。
\item \textbf{显存估算} 需要同时考虑参数、梯度、优化器状态、激活和 buffer，并区分哪些状态被 TP/PP/DP/ZeRO 切分。
\item \textbf{通信估算} 需要分开看 DP、TP、PP：DP 主要同步梯度，TP 高频层内 collective，PP 传递 stage 边界激活。
\item \textbf{计算效率} 受 GEMM shape、pipeline bubble、通信重叠、数据加载、checkpoint 和 straggler 影响，MFU 是重要观察指标。
\item \textbf{训练配置推导} 要满足 global batch、micro-batch、gradient accumulation、TP/PP/DP 整除关系和硬件拓扑约束。
\item \textbf{Distributed checkpoint} 是大规模训练可靠性的核心。checkpoint 需要保存模型、优化器、scheduler、RNG、consumed tokens、parallel metadata 和 shard layout。
\item \textbf{Checkpoint 策略} 需要在训练吞吐和故障回退成本之间平衡。频率太高会拖慢训练，频率太低会在故障后损失大量工作。
\end{itemize}

至此，我们已经建立了大模型训练最核心的并行框架：DP 解决数据扩展，TP 解决层内矩阵过大，PP 解决层数和整体模型过大，3D 并行则把三者组合为可扩展的训练系统。下一章将进入 ZeRO 优化器与 DeepSpeed，进一步分析如何在数据并行维度切分 optimizer state、gradient 和 parameter，从而显著降低每卡显存占用，并理解 ZeRO Stage 1、Stage 2、Stage 3 的通信与显存权衡。
